%% 上海海事大学《电子与电路学实验》课程报告通用模板
%% 编译方式：XeLaTeX
%% 说明：已删除具体实验内容，仅保留课程报告主要框架与通用排版

\documentclass[10pt,a4paper,twocolumn,twoside,UTF8]{ctexart}

%=========================================================
% 导言区
%=========================================================
\usepackage{geometry}
\geometry{left=2cm,right=2cm,top=2.5cm,bottom=3cm}

\usepackage{xeCJK}
\usepackage{amsmath,amssymb}
\usepackage{enumerate}
\usepackage{booktabs,multirow,array}
\usepackage{graphicx}
\usepackage{float}
\usepackage{setspace}
\usepackage{indentfirst}
\usepackage{titlesec}
\usepackage{fancyhdr}
\usepackage{ifthen}
\usepackage{datetime}
\usepackage[colorlinks=false,hidelinks]{hyperref}

\setlength{\parindent}{2em}
\setlength\columnsep{0.8cm}
\renewcommand{\arraystretch}{1.25}
\raggedbottom

%=========================================================
% 基本信息——每次实验只需修改这里
%=========================================================
\newcommand{\coursename}{电\quad 子\quad 与\quad 电\quad 路\quad 学\quad 实\quad 验}
\newcommand{\courseenglish}{ELECTRONICS AND CIRCUITS LABORATORY}
\newcommand{\experimenttitle}{实验题目}
\newcommand{\studentclass}{班级}
\newcommand{\studentname}{姓名}
\newcommand{\studentid}{学号}
\newcommand{\schoolname}{上海海事大学}
\newcommand{\collegename}{信息工程学院}
\newcommand{\schooladdress}{上海 浦东新区}
\newcommand{\postcode}{201306}

%=========================================================
% 标题样式
%=========================================================
\titleformat*{\section}{\large\bfseries\heiti}
\titleformat*{\subsection}{\normalsize\bfseries\heiti}
\titleformat{\paragraph}[block]{\normalsize\bfseries\heiti}{}{0em}{}
\titlespacing*{\paragraph}{0pt}{0.55em}{0.15em}

%=========================================================
% 页眉页脚
%=========================================================
\newboolean{first}
\setboolean{first}{true}
\pagestyle{fancy}

\fancypagestyle{maincontent}{
    \fancyhf{}
    \fancyhead[EL,OR]{\thepage}
    \fancyhead[EC]{\experimenttitle}
    \fancyhead[OC]{\coursename}
    \renewcommand\headrulewidth{0pt}
}

\fancypagestyle{firstpage}{
    \setboolean{first}{false}
    \fancyhf{}
    \fancyhead[L]{\the\year 年\the\month 月}
    \fancyhead[R]{\shortmonthname[\the\month], \the\year}
    \fancyhead[C]{
        \large{\coursename}\\
        \normalsize{\courseenglish}
    }
}

\newcommand{\makefirstpageheadrule}{
    \makebox[0pt][s]{\rule[0.6\baselineskip]{\headwidth}{0.3pt}}
    \makebox[0pt][s]{\protect\hspace{-0.34em}\rule[0.75\baselineskip]{\headwidth}{0.3pt}}
    \protect\vspace{-20pt}
}

\newcommand{\makeheadrule}{
    \makebox[0pt][l]{\rule[1\baselineskip]{\headwidth}{0.3pt}}
    \protect\vspace{-20pt}
}

\renewcommand{\headrule}{
    \ifthenelse{\boolean{first}}{\makeheadrule}{\makefirstpageheadrule}
}

%=========================================================
% 通用命令
%=========================================================
\newcommand{\blankline}[1][3cm]{\rule{#1}{0.4pt}}

% 图片占位框
\newcommand{\placeholderbox}[2]{%
    \begin{center}
    \fbox{\parbox[c][#2][c]{0.92\columnwidth}{\centering #1}}
    \end{center}
}

% 通用实验图片占位
\newcommand{\experimentimage}[2][实验图片]{%
    \begin{figure}[H]
        \centering
        \fbox{\parbox[c][#2][c]{0.92\columnwidth}{\centering\small #1}}
        \caption{#1}
    \end{figure}
}

% 通用过程材料图片占位
\newcommand{\processmaterialplaceholder}[1][实验过程记录]{%
    \begin{figure}[H]
        \centering
        \fbox{\parbox[c][4.6cm][c]{0.92\columnwidth}{\centering\small 在此插入实验现场照片}}
        \caption{#1}
    \end{figure}
}

%=========================================================
% 正文
%=========================================================
\begin{document}

\title{\LARGE\textbf{\experimenttitle}}
\author{\large\textit{\studentclass\quad \studentname\quad \studentid}\\
\normalsize{（\textit{\schoolname\ \collegename，\schooladdress\ }\postcode）}}
\date{}

\twocolumn[
\begin{@twocolumnfalse}
    \maketitle
    \renewcommand{\abstractname}{}
    \begin{abstract}
    \vspace{-3em}
    {\bf 摘{} 要：}
    {\small 在此填写本次实验的研究对象、主要实验方法、核心实验过程以及主要结果。摘要应简明概括实验内容，不宜大篇幅展开实验步骤。}
    \par
    \textbf{关键词}：关键词1；关键词2；关键词3；关键词4
    \vspace{1.5em}
    \end{abstract}
\end{@twocolumnfalse}
]

\thispagestyle{firstpage}
\pagestyle{maincontent}

%=========================================================
\section{实验目的}
% 按实验指导书填写，建议 2--5 条
\begin{enumerate}[(1)]
    \setlength{\itemsep}{0.15em}
    \setlength{\parsep}{0pt}
    \setlength{\topsep}{0.3em}
    \item 实验目的1。
    \item 实验目的2。
    \item 实验目的3。
\end{enumerate}

%=========================================================
\section{实验原理}
\indent 在此说明本实验涉及的主要理论、基本原理、公式关系或分析方法。

% 示例公式格式，可按需要保留或删除
\begin{equation}
    y=f(x)
\end{equation}

\indent 对公式中的主要物理量、变量及其含义进行必要说明。

%=========================================================
\section{实验仪器}
\indent 本实验所用仪器及元件如下表所示。
\vspace{-0.35em}

\begin{table}[H]
\centering
\footnotesize
\renewcommand{\arraystretch}{1.12}
\setlength{\tabcolsep}{3.2pt}
\begin{tabular}{p{0.32\columnwidth}p{0.13\columnwidth}p{0.41\columnwidth}}
\toprule
仪器或元件 & 数量 & 型号/参数 \\
\midrule
仪器/元件1 &  &  \\
仪器/元件2 &  &  \\
仪器/元件3 &  &  \\
仪器/元件4 &  &  \\
\bottomrule
\end{tabular}
\caption{实验仪器与元件}
\label{tab:devices}
\end{table}

\experimentimage[实验仪器与元件实物图]{4.0cm}

%=========================================================
\section{实验内容及结果}
\indent 在此简要说明本实验的总体实验内容、实验对象、实验条件及整体操作流程。

\experimentimage[实验原理图或实验电路图]{5.0cm}

%---------------------------------------------------------
\subsection{实验项目一}

\paragraph{1. 实验内容}
\indent 在此填写实验项目一的具体操作步骤、测量方法、变量设置及注意事项。

\paragraph{2. 实验数据与结果}
\indent 在此填写实验测量数据，并根据实验需要修改下列表格。

\begin{table}[H]
\centering
\small
\begin{tabular}{cccc}
\toprule
测量量 & 数据1 & 数据2 & 数据3 \\
\midrule
测量值 &  &  &  \\
\bottomrule
\end{tabular}
\caption{实验项目一原始数据}
\label{tab:experiment1}
\end{table}

\experimentimage[实验项目一现场记录]{4.0cm}

\par\indent 实验现象及结果说明：\par
\indent 在此描述实验现象，并结合理论关系、数据处理结果或误差情况对实验结果进行说明。

%---------------------------------------------------------
\subsection{实验项目二}

\paragraph{1. 实验内容}
\indent 在此填写实验项目二的实验步骤与操作方法。

\paragraph{2. 实验数据与结果}
\indent 在此填写实验项目二的数据及相关处理结果。

\begin{table}[H]
\centering
\small
\begin{tabular}{cccc}
\toprule
测量量 & 数据1 & 数据2 & 数据3 \\
\midrule
测量值 &  &  &  \\
\bottomrule
\end{tabular}
\caption{实验项目二原始数据}
\label{tab:experiment2}
\end{table}

\experimentimage[实验项目二现场记录]{4.0cm}

\par\indent 实验现象及结果说明：\par
\indent 在此分析实验现象、数据规律及其与理论预期之间的关系。

%---------------------------------------------------------
% 如实验包含更多子项目，可复制以下结构继续添加：
%
% \subsection{实验项目三}
% \paragraph{1. 实验内容}
% \indent ……
%
% \paragraph{2. 实验数据与结果}
% \indent ……

%=========================================================
\subsection{实验结果分析}
\indent 从实验数据、理论关系、实验现象和误差来源等方面，对本次实验结果进行综合分析。

\indent 可重点讨论：实验结果是否符合理论预期；不同实验条件之间有何变化规律；测量误差可能来自哪些环节；实验方法存在哪些局限。

%=========================================================
\section{实验结论}
\indent 在此简要总结本实验得到的主要结论。建议围绕实验目的逐项回应，不重复大段实验步骤。

%=========================================================
\section{思考与讨论}
\indent 在此填写实验指导书中的思考题、拓展问题或实验过程中发现的值得讨论的问题。

%=========================================================
\clearpage
\section{过程材料}
\indent 本部分集中展示实验过程中的接线、仪表读数、测量操作、实验台状态或其他过程性材料。根据实际需要复制或删除图片占位框。

\processmaterialplaceholder[实验过程记录1]
\processmaterialplaceholder[实验过程记录2]
\processmaterialplaceholder[实验过程记录3]

%=========================================================
% 如课程要求参考文献，可取消以下注释
%=========================================================
% \begin{thebibliography}{9}
% \bibitem{ref1} 作者. 文献名称[M/J]. 出版地: 出版社, 年份.
% \bibitem{ref2} 作者. 文献名称[J]. 期刊名称, 年份, 卷(期): 页码.
% \end{thebibliography}

\end{document}
